%%
% TUM tumarticle template — see https://gitlab.lrz.de/latex4ei/tum-templates
%%

\documentclass[
  english,        % define the document language (english, german)
  font=times,     % define main text font (helvet, times, palatino, libertine)
  twocolumn,      % use onecolumn or twocolumn layout
]{tumarticle}

\usepackage{cite}

% Add packages here:
\usepackage{booktabs}

% article metadata
\title{Group Assignment -- Business Process Prediction, Simulation, and Optimization}
\subtitle{A Data-Driven Simulation Model of the BPIC-17 Loan Application Process}

% TODO: fill in full names + emails of all four members
\author[email=sophie.opheiden@tum.de]{Sophie~Opheiden}
\author[email=todo@tum.de]{Mario~(TODO)}
\author[email=todo@tum.de]{Daniel~(TODO)}
\author[email=todo@tum.de]{Johannes~(TODO)}

\date\today

\begin{document}

\maketitle

\begin{abstract}
% TODO (all): 3–4 sentences once Part II results are in. Suggested skeleton:
% We present a data-driven discrete-event simulation model of the BPIC-17
% loan application process and use it as a testbed for resource allocation
% policies. All simulation parameters are mined from the real event log; the
% model is validated with a second-pass comparison against the original log
% and reproduces the real mean case duration within 2.5\%. On top of the
% simulator, we compare [n] allocation policies on cycle time, resource
% occupation, and fairness, and answer [management question].
tbd.
\end{abstract}
%

\section{Introduction}
Business process simulation makes it possible to answer \emph{what-if}
questions -- how would cycle times change with fewer employees, different
working hours, or a different work distribution -- without touching the real
process. The reliability of such answers, however, is bounded by how
faithfully the simulation reproduces reality. Instead of modeling by hand, we
therefore follow the data-driven approach of Rozinat et
al.~\cite{rozinat2009discovering}: every model component -- control-flow,
branching behavior, arrival process, processing times, and the resource
perspective -- is mined from a real event log, and the resulting simulator is
validated by comparing the statistics of its simulated logs against the
original log (the ``second pass'' method).

As data source we use the BPI Challenge 2017 event
log~\cite{vandongen2017bpic}, covering 31{,}509 loan applications
(1.2M events, 26 activities, 149 resources) of a Dutch financial institution
over 13 months. The process combines milestone events for the application
(\texttt{A\_}\dots) and offer (\texttt{O\_}\dots) lifecycles with manually
executed work items (\texttt{W\_}\dots), and contains substantial loop
behavior (repeated validation and offer rounds), concurrency, and long
customer-response waiting times -- properties that, as we show, break naive
simulation approaches and require explicit design decisions.

This report is structured along the two assignment parts:
Section~\ref{sec:simulation} describes the simulation model, one subsection
per component, with the responsible group member named at the start of each
subsection; Section~\ref{sec:optimization} covers the resource allocation
optimizations built on top of it. Design decisions are justified empirically
throughout: every number cited in this report is reproducible from a metrics
file in the repository's \texttt{output/} directory.

\section{Simulation Model}
\label{sec:simulation}

\subsection{Simulation Engine Core (Sophie)}
The engine implements discrete event simulation with a single global event
queue, as required by the assignment: a min-heap ordered by
$(\text{timestamp}, \text{priority})$, where the integer priority breaks
ties between simultaneous events deterministically (e.g.\ a resource is
released before the next task is dispatched). Event types cover the case
lifecycle (\texttt{CASE\_ARRIVAL}, \texttt{CASE\_COMPLETE}), the activity
lifecycle (\texttt{ACTIVITY\_START}, \texttt{ACTIVITY\_COMPLETE}), and
resource state changes. The engine itself contains no domain logic -- it is
a dispatcher: components (arrivals, process, resources) are plain classes
that register handlers per event type and are called in registration order.
This design decision let all four members develop their components
independently against a stable interface, without ever modifying the engine
or each other's code. A built-in logger records every activity event in
pm4py column convention and exports CSV; case ids that reach
\texttt{CASE\_COMPLETE} are written alongside, so that every evaluation can
filter out horizon-truncated cases (unfiltered logs would bias cycle time
and case duration downwards). Simulation time is a float in seconds --
the natural base unit for datetime arithmetic -- anchored to a real start
date at logging time. All stochastic components draw from seeded random
number generators, and set iterations that would otherwise depend on memory
layout are explicitly sorted, so a fixed seed reproduces a byte-identical
event log -- the basis for all before/after comparisons in this report.

\subsection{Case Arrivals (Mario)}
\textbf{Basic.} Case arrivals follow a single LogNormal inter-arrival
distribution fitted on the real log ($s=1.4543$, scale $363.10$s, mean
inter-arrival $1002.23$s $\approx 86$ cases/day, $\sigma=32$/day). This is
time-independent by construction: it can match the mean rate but cannot
represent that arrivals are not uniform in time.

\textbf{Advanced.} The real process is locally Poisson but with a rate
that varies strongly by time of day and weekday (night $\sim\!0.6$
arrivals/h vs.\ $\sim\!7.6$/h in the 12--18h core; Monday
$\approx 3\times$ Sunday), which a static distribution cannot capture in
principle. We replace it with an intensity-free temporal point process: a
small Mixture-Density Network (a mixture of LogNormals) predicts the
inter-arrival distribution conditioned on cyclical time-of-day/weekday/
season features of the current simulation clock, trained offline
(\texttt{train\_arrival\_mdn.py}) and evaluated at runtime as a pure
NumPy forward pass (no PyTorch dependency at simulation time). It is a
drop-in replacement for the parametric component (same
\texttt{bootstrap}/\texttt{on\_arrival} interface), toggled by
\texttt{USE\_MDN\_ARRIVALS} in \texttt{main.py} (default off).

\textbf{Validation status.} \texttt{scripts/eval\_arrivals.py} implements
the intended KPI comparison -- real vs.\ parametric vs.\ MDN arrivals on
hour-of-day MAE, weekday MAE, and an inter-arrival KS statistic, deliberately
refusing to substitute the simulated log for ground truth -- but requires
the raw BPIC-17 log, which is gitignored and was not available on the
machine this was implemented on. This is the one open item on this
subsection: whoever has the raw log should run it once
(\texttt{BPIC17\_LOG=\ldots{} python scripts/eval\_arrivals.py}) to produce
\texttt{output/arrival\_model\_eval.md} before the report is finalized, and
the team should jointly decide whether to flip the default based on the
result -- it changes every teammate's simulated logs.

\subsection{Processing Times (Daniel)}
% TODO (Daniel): fitted distributions (Basic 1), contextual GBR point model
% (Basic 2), quantile-regression distribution model (Advanced I); temporal
% 80/20 split (random split overstated log-$R^2$ 0.36 vs.\ 0.12!), MAE/MSE,
% pinball loss + 90\%/50\% interval coverage.
% Evidence: output/models/processing_time_metrics.json
tbd.

\subsection{Process Model (Sophie)}
\label{sec:processmodel}
The \emph{basic} variant routes cases through a flat next-activity graph with
empirical transition probabilities. It has no structural notion of legality:
$0\%$ of its simulated traces replay on a reference model of the process
(average token-replay fitness $0.64$), i.e., it freely generates behavior
that cannot occur in reality. The \emph{advanced} variant therefore loads a
BPMN 2.0 model, converts it to a Petri net using
pm4py~\cite{berti2019pm4py}, and enforces control-flow through the firing
rules of the net: each case holds its own marking, and the set of legal next
activities is the set of enabled visible transitions. Silent
($\tau$)~transitions from gateway logic are handled by computing the full
$\tau$-closure of the current marking -- a greedy strategy that fires
$\tau$-transitions only when no visible transition is enabled would make
whole branches structurally unreachable, because n-way splits are encoded as
nested binary gateways whose named branch and $\tau$-sibling are enabled
simultaneously. Following the lecture's advice to guard against structural
design errors, we verified the converted net with pm4py's Woflan
implementation: it is a sound workflow net (live, bounded, no dead
transitions) -- the loop behavior analyzed in Section~\ref{sec:branching}
is therefore legitimate model behavior, not a modeling defect.

\textbf{Design decision: model source.} We compared our manually modeled
(Signavio) BPMN against models discovered from the log with the Inductive
Miner. The discovered model (noise threshold $0.2$) replays the real log
only marginally better ($73.6\%$ vs.\ $68.8\%$ fully fitting traces,
average fitness $0.993$ vs.\ $0.976$), but performs clearly worse when
simulating: precision drops from $0.78$ to $0.49$ and the branching
divergence (mean total variation distance, TVD) doubles from $0.24$ to
$0.57$, while a termination failure we describe in
Section~\ref{sec:branching} occurs identically under both models
($2\%$ completion). We concluded that (i)~the manual model is the better
simulation basis, and (ii)~non-termination is a property of the branching
mechanism, not of the model source -- a conclusion we could only draw
because both options were evaluated with the same KPI suite.
Figure~\ref{fig:bpmn} shows the model.

\begin{figure*}[t]
    \centering
    % TODO: export the Signavio model as PNG/PDF into report/figures/
    \IfFileExists{figures/bpic17_signavio_bpmn.png}
      {\includegraphics[width=\textwidth]{figures/bpic17_signavio_bpmn.png}}
      {\framebox[\textwidth]{\rule{0pt}{3cm}TODO: export Signavio BPMN to figures/bpic17\_signavio\_bpmn.png}}
    \caption{Manually modeled BPMN 2.0 model of the BPIC-17 loan application
    process (Signavio). Highlighted: the offer/validation loops whose exit
    behavior required visit-dependent branching (cf.\
    Fig.~\ref{fig:loopfragment}).}
    \label{fig:bpmn}
\end{figure*}

\subsection{Branching Decisions (Sophie)}
\label{sec:branching}
\textbf{Basic.} Branching at each decision point (a marking with several
enabled visible transitions) initially used memoryless next-activity
probabilities mined from trace bigrams, renormalized over the legal
candidates. This configuration exposed the central challenge of our
simulator: only $2\%$ of cases terminated within the simulation horizon, at
an average of $74.9$ events per case (real: $15.1$). Three measured root
causes: (i)~trace bigrams mix concurrency interleavings into the branching
estimate -- the most frequent real successor \emph{event} of
\texttt{A\_Validating} is \texttt{O\_Returned} ($93.6\%$), an activity of a
concurrent branch that is not even enabled at that decision point, so
probabilities were renormalized over the wrong candidates; (ii)~memoryless
probabilities understate loop exits -- in the real log the exit probability
escalates with each loop round (e.g.\ \texttt{W\_Validate application}
$\rightarrow$ \texttt{O\_Cancelled}: $3.5\% \rightarrow 14\% \rightarrow
27\%$ across visits); (iii)~termination itself is a \emph{decision}: at many
markings the final marking is only $\tau$-reachable while a visible
loop-back label remains enabled (the \texttt{[O\_Cancelled]} frontier was
evaluated $15.7\times$ per case with nothing else to choose), so checking
marking equality alone can never stop such cases.

\textbf{Fix (three parts).} First, branching probabilities are mined
\emph{at the decision points} by replaying the real log on the very Petri
net used for enforcement, conditioned on how often the case has already
visited that decision point. Second, wherever a real trace ends although
visible transitions remain enabled, an explicit \texttt{END} choice is
recorded (at \texttt{[O\_Cancelled]}, real cases stop with $80\%$ per
visit); the simulation offers \texttt{END} only with its mined probability
-- ending must be evidenced, never a residual. Third, a domain rule ends a
case once one of the outcome milestones \texttt{A\_Pending},
\texttt{A\_Denied}, \texttt{A\_Cancelled} has fired (in the real log
exactly one of them occurs per case: $0.55+0.33+0.12=1.00$).
Table~\ref{tab:ablation} isolates the contribution of each part: neither
the terminal rule alone nor visit-conditioning alone suffices; combined,
completion rises from $2\%$ to $71\%$, the case-length error falls from
$3.96$ to $0.15$, branching TVD from $0.24$ to $0.11$, and the simulation
reproduces 10 of the 20 most frequent real trace variants -- previously
zero. Figure~\ref{fig:loopfragment} visualizes the mined loop semantics.

\begin{table}[h]
    \caption{Branching ablation (advanced model, seed 42, 30-day horizon;
    KPIs vs.\ the real log). \emph{term} = terminal-outcome rule,
    \emph{visit} = decision-point probabilities with visit buckets +
    mined \texttt{END}. Basic = flat graph for reference.}
    \centering
    \renewcommand{\arraystretch}{1.2}
    \setlength{\tabcolsep}{3.5pt}
    \small
    \begin{tabular}{lrrrrr}
        \toprule
        \textbf{KPI} & \textbf{Basic} & \textbf{probs} & \textbf{+term} & \textbf{visit} & \textbf{rules} \\
        \midrule
        Fitness (fully fitting) & $0\,\%$ & $100\,\%$ & $0.6\,\%$ & $45.6\,\%$ & $48.2\,\%$ \\
        Fitness (avg.\ trace)   & $0.64$ & $1.00$ & $0.95$ & $0.97$ & $0.97$ \\
        Precision               & $0.72$ & $0.78$ & $0.72$ & $0.69$ & $0.47$ \\
        Branching mean TVD      & $0.41$ & $0.24$ & $0.38$ & $\mathbf{0.11}$ & $0.15$ \\
        Top-20 variants repr.   & $0$ & $0$ & $0$ & $\mathbf{10}$ & $\mathbf{10}$ \\
        Case length (real 15.1) & $12.3$ & $74.9$ & $16.9$ & $12.9$ & $\mathbf{15.2}$ \\
        Completion rate         & $94\,\%$ & $2\,\%$ & $37\,\%$ & $71\,\%$ & $\mathbf{73\,\%}$ \\
        \bottomrule
    \end{tabular}
    \label{tab:ablation}
\end{table}

\begin{figure}[t]
    \centering
    % TODO (Sophie, Signavio): model the validation/offer loop fragment with
    % the mined semantics — see chat/report notes for the concrete content.
    \IfFileExists{figures/loop_fragment_signavio.png}
      {\includegraphics[width=\columnwidth]{figures/loop_fragment_signavio.png}}
      {\framebox[\columnwidth]{\rule{0pt}{3cm}TODO: figures/loop\_fragment\_signavio.png}}
    \caption{Validation/offer loop fragment (Signavio) with the mined
    visit-dependent semantics: XOR gateways annotated with
    $P(\text{branch} \mid \text{visit } k)$, the loop-back edge whose exit
    probability escalates ($3.5\%\!\rightarrow\!14\%\!\rightarrow\!27\%$),
    the data-driven \texttt{END} decision at \texttt{[O\_Cancelled]}
    ($80\%$ per visit), and the terminal outcome milestones.}
    \label{fig:loopfragment}
\end{figure}

\textbf{Advanced I: attribute-based decisions.} Case attributes are sampled
at spawn from distributions learned from the log
(\texttt{LoanGoal} and \texttt{RequestedAmount} conditional on
\texttt{ApplicationType}), and executing \texttt{O\_Create Offer} sets the
offer attributes at runtime. Per decision point, a decision tree predicts
the branch from these attributes (\emph{rules} column in
Table~\ref{tab:ablation}); the training data was identified by replaying the
log on the process model, and outcome-encoding attributes
(\texttt{Accepted}, \texttt{Selected}) were excluded to prevent label
leakage. Evaluated with a temporal 80/20 split -- training on older, testing on the
most recent cases, since a random split would leak process drift -- and the
metrics from the lecture (precision, recall, one-vs-rest ROC-AUC next to
accuracy and a majority-class baseline; no separate validation slice, as all
hyperparameters were fixed a priori rather than tuned), 16 decision points
obtained models: one is
perfectly separable (AUC $1.0$), while several collapse to the majority
class -- e.g.\ accuracy $0.89$ at macro-precision $0.22$ -- which accuracy
alone would have hidden. \emph{rules} matches \emph{visit} on variant
reproduction and is even closer in case length (error $0.004$); we kept
\emph{visit} as default for its better TVD and precision.

\textbf{Validation and honest limitations.} The fitness trade-off in
Table~\ref{tab:ablation} (fully fitting $100\%\!\rightarrow\!46\%$ at
average fitness $0.97$) arises because real -- and now also simulated --
traces end before the net's final marking; it documents a limitation of the
model, not of the simulator. Because a 30-day horizon only lets fast cases
finish, we additionally ran a drain experiment (arrivals for 30 days,
simulation drained to day 180): completion reaches $99.2\%$ and the mean
case duration matches the real process at $21.3$ vs.\ $21.8$ days --
a relative error of $2.5\%$; the seemingly large duration errors of
horizon-bounded runs are survivorship bias. The drained population also
reveals the remaining gap: long-running cases loop more than in reality
(mean length 46 events), because beyond the fifth visit our buckets keep
the exit probability stationary.

\subsection{Resource Availability, Permissions, and Allocation (Johannes)}
% TODO (Johannes): interval-based availability (Basic 1.6) / mined calendars,
% permission map from the log (1.7), random allocation (1.8) as the Part II
% baseline; resource contention rework (activities wait for a resource).
tbd.

\section{Optimization}
\label{sec:optimization}

\subsection{Experimental Setup and Metrics (Mario)}
\label{sec:optsetup}
\texttt{scripts/run\_experiments.py} orchestrates every policy comparison
in this section: a policy $\times$ seed $\times$ scenario grid, one
simulation run per cell, evaluated with the three lecture-slide-21
metrics -- average cycle time (from case arrival, not first activity
start), average resource occupation (relative to the Section~1.6
availability windows, not wall-clock span), and (weighted) resource
fairness -- computed by \texttt{scripts/opt\_metrics.py}, plus four
self-designed metrics justified against process-specific failure modes
the slide-21 trio cannot see on its own: time-to-first-offer and
time-to-decision (customer-facing, from the \texttt{A\_}/\texttt{O\_}
milestone events), handover rate (share of consecutive case steps that
switch resource -- familiarity/context-switch cost), and rolling
workload balance (per-day resource-occupation std, catching ``fair on
average, unfair in bursts'' that a single whole-horizon fairness number
cannot). Scenarios beyond
the default (\emph{normal}) are \emph{peak} ($+30\%$ arrival rate) and
\emph{outage} (20\% of resources removed, seeded per run), giving the
infrastructure the management question and RQ4 need without bespoke code.

\textbf{Common Random Numbers.} A policy change alters resource-allocation
\emph{order}, which -- because every stochastic draw in
\texttt{ProcessComponent} previously came from one shared RNG consumed in
event-dispatch order -- silently changed every subsequent branching and
duration draw too: two policies under ``the same seed'' were really
running different case trajectories, not a controlled comparison (we
measured this directly for Piled Execution: see
\texttt{output/piled\_execution\_eval.md}). We therefore seed each
branching/duration decision independently from
(base seed, case id, activity, visit count), gated behind \texttt{crn=True}
(default off, to keep every prior evidence log bit-for-bit reproducible).
We verified directly that this achieves its goal: under CRN, two runs that
differ only in allocation policy produce an \emph{identical} branching
sequence for every case, up to whichever run's allocation timing causes it
to hit the horizon with fewer completed activities first
(\texttt{scripts/test\_crn.py}).

\textbf{Caveat.} The results in
Section~\ref{sec:kbatching} use only 5 seeds (the runner itself flags
this: the 95\% CI half-width exceeds the roadmap's 5\%-of-mean target on
every policy at this seed count) -- treat the ranking as indicative, not
conclusive, until re-run with $\geq 10$ seeds.

\subsection{Simple Policies and k-Batching (Mario: k-Batching)}
\label{sec:kbatching}
The Section~1.8 baseline is R-RMA (uniform random pick among qualified,
available resources, Russell et al.~\cite{russell2005workflow} pattern~15)
behind a minimal \texttt{AllocationPolicy} seam
(\texttt{simulation/policies.py}) that the simple heuristics (R-RRA, R-SHQ)
and the two advanced policies plug into without touching
\texttt{ResourceComponent} again.

\textbf{k-Batching (Zeng \& Zhao).} Rather than allocating a work item the
instant it is enabled, k-Batching lets items queue and releases them in
batches of $k$, solved as a parallel-machines assignment problem: batch
items $\times$ currently free-and-on-shift resources, cost = expected
processing time from a shared \texttt{expected\_duration()} model (the
same trained GBR point estimator Section~1.3 built, wrapped in
\texttt{simulation/expected\_duration.py} so both this and the Park~\&
Song policy consume one implementation), solved with
\texttt{scipy.optimize.linear\_sum\_assignment}. A flush triggers once $k$
items are waiting \emph{or} the oldest waiting item has waited past a 4h
safety valve -- without it, a queue that never reaches $k$ under low load
would starve indefinitely, the known weakness the lecture attributes to
this heuristic.

Table~\ref{tab:kbatch} sweeps $k \in \{1,2,5,10,20\}$ against the R-RMA
baseline (5 seeds, 30-day horizon, capacity 3, calendar availability, CRN
on). Every $k$ increases mean cycle time relative to immediate allocation
(statistically significant at $k\!\in\!\{1,2,10,20\}$, paired $t$-test
$p<0.05$; $k\!=\!5$ borderline, $p=0.059$) -- exactly the idle-time cost
the lecture predicts: k-Batching's requester always waits for a flush
trigger, even when a resource happens to be free right now, so it cannot
beat an immediate-allocation baseline on cycle time by construction, only
on assignment \emph{quality} once it does flush. That quality effect is
visible in the trend: cycle time falls monotonically from $k\!=\!1$ to
$k\!=\!10$ ($10.49\text{d} \rightarrow 9.01\text{d}$) as larger batches
let the assignment solver place more items well per flush, before rising
again at $k\!=\!20$ once queueing delay dominates. Occupation shows the
same shape in reverse ($2.63 \rightarrow 4.18 \rightarrow 4.56$,
approaching the baseline's $5.19$): larger batches keep more resources
usefully busy per flush. This is an honest negative-leaning result rather
than a clean win, and we report it as such: k-Batching's real argument
here is assignment quality under heterogeneous resource-activity
durations, not automatic throughput -- exactly the trade-off that
motivates the strategic-idling advanced policies below (Park \& Song
plan ahead using predicted next tasks instead of waiting for a batch to
fill; see the corresponding subsection).

% TODO (Daniel): \subsection{Advanced Policy I: Park \& Song (Daniel)} --
%   next-task prediction (reuse the Section 1.3 duration models) + strategic
%   idling, assignment/min-cost-max-flow.
% TODO (Johannes): \subsection{Advanced Policy II: Kunkler \& Rinderle-Ma
%   (Johannes)} -- dummy-resource-cost formulation for uncertain
%   availability, builds on the Section 1.6 calendars.
% TODO (Sophie): \subsection{Management Question (Sophie)} -- "fire two
%   employees": criticality analysis (occupation, exclusive permissions,
%   historical performance) + leave-two-out simulations using
%   ResourceComponent(excluded_resources=...) / --scenario outage
%   infrastructure from scripts/run_experiments.py.

\begin{table}[h]
    \caption{k-Batching sweep vs.\ R-RMA baseline (5 seeds, 30-day
    horizon, seed-42-derived; CIs wide at this seed count, see
    Section~\ref{sec:optsetup}). Cycle time in days; occupation and
    fairness as in slide 21 (occupation $>\!1$: capacity-3 parallelism,
    not an error, see \texttt{opt\_metrics.py}).}
    \centering
    \renewcommand{\arraystretch}{1.2}
    \setlength{\tabcolsep}{4pt}
    \small
    \begin{tabular}{lrrr}
        \toprule
        \textbf{Policy} & \textbf{Cycle time (d)} & \textbf{Occupation} & \textbf{Fairness} \\
        \midrule
        R-RMA (baseline) & $7.93$          & $5.19$ & $1.15$ \\
        k=1              & $10.49$         & $2.63$ & $1.12$ \\
        k=2              & $9.82$          & $3.26$ & $1.08$ \\
        k=5              & $9.53$          & $3.86$ & $1.00$ \\
        k=10             & $\mathbf{9.01}$ & $4.18$ & $0.97$ \\
        k=20             & $9.30$          & $4.56$ & $1.02$ \\
        \bottomrule
    \end{tabular}
    \label{tab:kbatch}
\end{table}

\section{Use of AI Tools}
% TODO (all): adapt to what each member actually did. Draft:
We used AI coding assistants (Claude Code, Anthropic) as pair-programming
support during this project: for implementing and refactoring simulation
components and evaluation scripts, for debugging, for the diagnostic
analyses that located the termination and concurrency problems described in
Section~\ref{sec:branching}, and for drafting parts of this report. All
design decisions were made by the team, and every AI-assisted change was
reviewed by the responsible member and validated against the KPI suite
before being adopted.

\section{Conclusion}
% TODO (all, ~0.3 col): once Part II is done. One paragraph: validated
% simulator (headline: 2.5% duration error, 10/20 variants), what the policy
% comparison showed, answer to the management question, limitations/outlook.
tbd.

%
\bibliographystyle{plain}
\bibliography{references}

\end{document}
